\section{从时空间隔到 Poincaré 代数}

取坐标 \(x^\mu=(ct,x^1,x^2,x^3)\)，Minkowski 度规
\(\eta=\operatorname{diag}(1,-1,-1,-1)\)，指标用 \(\eta\) 升降。
两事件的间隔 \(\eta_{\mu\nu}\Delta x^\mu\Delta x^\nu\) 在惯性系变换下不变。线性变换 \(\Lambda\) 满足
\(\Lambda^\top\eta\Lambda=\eta\)；若还保持时间方向及空间取向，称为正规正时 Lorentz 变换，这些矩阵组成的群记为 \(SO^+(1,3)\)。

把 \(\Lambda(\varepsilon)=I+\varepsilon A+O(\varepsilon^2)\)
代入间隔不变条件，比较 \(\varepsilon\) 的系数得到
\(A^\top\eta+\eta A=0\)。若定义
\(A_{\mu\nu}=\eta_{\mu\rho}A^\rho{}_\nu\)，这就是
\(A_{\mu\nu}=-A_{\nu\mu}\)：四维共有六个独立无穷小参数，
三个空间转动与三个沿空间轴的洛伦兹推动。例如沿 \(x^1\) 的推动取
\[
x^{\prime0}=x^0\cosh\xi+x^1\sinh\xi,\qquad
x^{\prime1}=x^0\sinh\xi+x^1\cosh\xi.
\]
直接展开两个平方并用 \(\cosh^2\xi-\sinh^2\xi=1\)，
就验证 \((x^{\prime0})^2-(x^{\prime1})^2
=(x^0)^2-(x^1)^2\)。在 \(\xi=0\) 求导给
\(x^1\partial_0+x^0\partial_1=J_{01}\)，
说明下面的生成元确实来自一个具体的连续变换。
参数 \(\xi\) 称为快度；把一条 \(x^1=0\) 的世界线代入变换式，得到变换后 \(x^{\prime1}/x^{\prime0}=\tanh\xi\)，对应速度大小 \(v=c|\tanh\xi|<c\)。双曲函数不是装饰：它同时保证间隔不变和速度不会越过 \(c\)。

\begin{definition}[Poincaré 群]
Poincaré 变换是仿射映射 \((a,\Lambda):x\mapsto\Lambda x+a\)，其中 \(a\in\mathbb R^{1,3}\)、\(\Lambda\in SO^+(1,3)\)。群乘法由映射复合给出
\[
(a,\Lambda)(b,\Gamma)=(a+\Lambda b,\Lambda\Gamma).
\]
这称为平移群与 Lorentz 群的半直积，因为后一因子会改变前一因子的向量。
\end{definition}

直接比较 \((a,I)(0,\Lambda)\) 与 \((0,\Lambda)(a,I)\)：所得平移分别是 \(a\) 与 \(\Lambda a\)，一般不同。这是平移生成元与 Lorentz 生成元不对易的原因。

\begin{definition}[Lie 代数与生成元]
光滑群的单位元处切空间称其 Lie 代数。取群中经过单位元的曲线，令它们在函数上的无穷小作用为向量场；括号定义为算子交换子 \([A,B]=AB-BA\)。平移方向取
\(P_\mu=\partial_\mu\)，Lorentz 方向取
\(J_{\mu\nu}=x_\mu\partial_\nu-x_\nu\partial_\mu=-J_{\nu\mu}\)。
这些是数学上的微分算子生成元，不含 \(\ii\hbar\)。
\end{definition}

用 \(\partial_\rho x_\mu=\eta_{\mu\rho}\) 逐项计算，对任意光滑函数 \(f\) 有
\[
[J_{\mu\nu},P_\rho]f
=J_{\mu\nu}(\partial_\rho f)
-\partial_\rho(J_{\mu\nu}f)
=(\eta_{\nu\rho}P_\mu-\eta_{\mu\rho}P_\nu)f.
\]
平移偏导彼此交换，故 \([P_\mu,P_\nu]=0\)。同理，展开 \(J_{\mu\nu}\) 作用于
\(x_\rho\partial_\sigma-x_\sigma\partial_\rho\) 并消去二阶导数，得到
\[
[J_{\mu\nu},J_{\rho\sigma}]
=\eta_{\nu\rho}J_{\mu\sigma}-\eta_{\mu\rho}J_{\nu\sigma}
-\eta_{\nu\sigma}J_{\mu\rho}+\eta_{\mu\sigma}J_{\nu\rho}.
\]
由具体微分算子得出的括号，在任何可微表示中都保留：若
\(\rho(g_1g_2)=\rho(g_1)\rho(g_2)\)，对群交换子曲线分别求两次导数便有
\(d\rho([A,B])=[d\rho(A),d\rho(B)]\)。这说明态空间生成元的代数结构来自群乘法；选用量子力学的 Hermite 生成元时，额外的 \(\ii\hbar\) 因子取决于指数变换约定。

群究竟怎样作用在场上也要指定。标量场可取主动表示
\[
(U(a,\Lambda)\phi)(x)=\phi(\Lambda^{-1}(x-a)).
\]
将 \((b,\Gamma)\) 的作用先代入，再作用 \((a,\Lambda)\)，
逆变换按相反顺序复合，正好得到
\(U(a,\Lambda)U(b,\Gamma)=U(a+\Lambda b,\Lambda\Gamma)\)。
矢量场除了自变量的逆变换，还要乘一个 \(\Lambda^\mu{}_\nu\)；
因此知道群元怎样相乘，尚不足以知道一个具体物理场怎样变换，
必须另外给它的表示。

一个可先验证的群不变量是平移生成元的二次式
\(P^2=\eta^{\mu\nu}P_\mu P_\nu\)。
因所有 \(P_\mu\) 彼此对易，\(P^2\) 与各 \(P_\rho\) 对易。对 \(J_{\alpha\beta}\) 用乘积的
交换子规则 \([J,AB]=[J,A]B+A[J,B]\)，再代入上面的
\([J_{\alpha\beta},P_\mu]\)，两项在对称缩并后相消，故
\([J_{\alpha\beta},P^2]=0\)。在单粒子表示里，相应 Hermite
四动量的平方取质量壳值 \(m^2c^2\)；这才说明“质量”
为何是 Poincaré 对称下不会随惯性系改变的量。

一个容易混淆的反例是把所有 Lorentz 变换当作连续对称：空间反演与时间反演不在 \(SO^+(1,3)\) 的单位元连通分支内，不能由上述无穷小 \(J_{\mu\nu}\) 连续生成。

\begin{exercise}
从 \(J_{12}=x_1\partial_2-x_2\partial_1\) 与 \(P_1=\partial_1\) 直接计算交换子，并用群乘法解释结果。再检查 \(P^\mu P_\mu\) 与全部 \(P_\rho,J_{\mu\nu}\) 对易。
\end{exercise}
